% This must be in the first 5 lines to tell arXiv to use pdfLaTeX, which is strongly recommended.
\pdfoutput=1
% In particular, the hyperref package requires pdfLaTeX in order to break URLs across lines.

\documentclass[11pt]{article}

% Remove the "review" option to generate the final version.
\usepackage{ACL2023}

% Standard package includes
\usepackage{times}
\usepackage{latexsym}
\usepackage{amsmath,amsthm,amssymb}
\usepackage{hyperref}

% For proper rendering and hyphenation of words containing Latin characters (including in bib files)
\usepackage[T1]{fontenc}
% For Vietnamese characters
% \usepackage[T5]{fontenc}
% See https://www.latex-project.org/help/documentation/encguide.pdf for other character sets

% This assumes your files are encoded as UTF8
\usepackage[utf8]{inputenc}

% This is not strictly necessary, and may be commented out.
% However, it will improve the layout of the manuscript,
% and will typically save some space.
\usepackage{microtype}

% This is also not strictly necessary, and may be commented out.
% However, it will improve the aesthetics of text in
% the typewriter font.
\usepackage{inconsolata}


% If the title and author information does not fit in the area allocated, uncomment the following
%
%\setlength\titlebox{5cm}
%
% and set <dim> to something 5cm or larger.
\title{Evaluating Qwen3 Reasoning Capabilities Outside the Thinking Mode}

\author{
  % TODO add you names here
  \textbf{Łukasz Smoliński}\\
  \texttt{ls438719@students.mimuw.edu.pl}\\
  \And
  \textbf{Marek Lisowski}\\
  \texttt{ml438584@students.mimuw.edu.pl}\\
  \AND
  \textbf{Jakub Nowacki}\\
  \texttt{jn431513@students.mimuw.edu.pl}
  \And
  \textbf{Yun Chia Lee}\\
  \texttt{y.lee15@student.uw.edu.pl}
}
\date{\vspace*{-0.2cm}}


\begin{document}
\maketitle

\begin{abstract}
We investigate whether Qwen3 models, trained to reason during post-training
with the special \texttt{<think>} token, can exhibit similar reasoning
behavior when prompted with natural language cues like \texttt{"Wait"}.
While the Qwen3 thinking mode is expected to perform better, we aim to
quantify how well Qwen3 models reason using regular language cues under
various thinking budgets. Additionally, we compare Qwen3's performance with
Qwen2.5 to assess the generalizability of reasoning capabilities elicited
through natural language prompts. This will help determine whether
post-training in Qwen3 improves general-purpose reasoning without
explicit fine-tuning.
\end{abstract}

\section{Introduction}
Large language models (LLMs) have demonstrated remarkable capabilities across
various natural language processing tasks. A critical aspect of their
performance lies in their ability to perform complex reasoning. Traditionally,
enhancing reasoning in LLMs has involved techniques like Chain-of-Thought
(CoT) prompting, where models are guided to generate intermediate reasoning
steps before arriving at a final answer \citep{cot}. This approach
has shown significant improvements in tasks requiring multi-step reasoning.

Recent advancements have introduced models like Qwen3, which incorporate a
dedicated thinking mode activated by special \texttt{<think>} and
\texttt{</think>} tokens. This mode, integrated during post-training, allows the
model to perform complex, multi-step reasoning tasks more effectively
\citep{qwen3}.

However, it remains unclear whether such specialized reasoning abilities can
be elicited without relying on these special tokens. Specifically, we aim to
investigate if natural language cues, such as the word \texttt{"Wait"}, can
trigger similar reasoning processes in Qwen3 models. This exploration is
motivated by the desire to understand the extent to which reasoning
capabilities are inherently built into the model versus being reliant on
\texttt{<think>} tokens.

We propose to evaluate the Qwen3-1.7B model under different prompting
strategies, including:
\begin{itemize}
    \item \textbf{Dedicated thinking mode}, where the model is explicitly
    instructed to engage in structured reasoning with \texttt{<think>} tokens.
    \item \textbf{Natural language cues}, where reasoning is encouraged
    implicitly through natural language prompting.
    \item \textbf{Baseline}, where the model is prompted to answer
    directly without any explicit encouragement to reason.
\end{itemize}
Additionally, we will assess how varying the thinking budget --- the
number of tokens allocated for reasoning --- impacts the model's performance
across these different prompting methods.

Furthermore, to evaluate the generalizability of reasoning capabilities
elicited through natural language prompts, in our experiments we will
include Qwen2.5-1.5B \citep{qwen2.5}, which has previously shown improvements
under the budget forcing technique after fine-tuning \citep{s1}, but lacks
Qwen3's reasoning enhancements. This comparative analysis will help determine
whether the observed behaviors are unique to Qwen3.

We will conduct experiments using the SVAMP \citep{svamp} dataset, which tests
mathematical and reasoning abilities of language models. By analyzing
performance, we aim to gain insights into the effectiveness of different
prompting strategies and the inherent reasoning capabilities of the models
under study.

\section{Related Work}
Recent work by \citet{s1} introduced the idea of budget forcing, where models
are encouraged to continue reasoning by appending the word \texttt{"Wait"}
multiple times during inference. Their experiments with the Qwen2.5-32B-Instruct
model, fine-tuned on a small dataset of $1000$ reasoning examples, showed
that this approach could improve performance on benchmarks consisting of mathematical problems.

The Qwen3 series \citep{qwen3} shifts the paradigm by embedding reasoning
capabilities directly into the model during post-training. This is done
via dedicated \texttt{<think>} tokens, allowing the model to enter a
reasoning phase and produce structured thought processes. The extent to
which this behavior generalizes beyond the \texttt{<think>} block,
however, remains unclear.

Our study aims to bridge this gap by testing whether Qwen3's special post-training
allows for reasoning even outside the thinking mode and without any
additional fine-tuning, and how it compares to Qwen2.5, which lacks
this special feature by default.

\bibliography{refs}
\bibliographystyle{acl_natbib}
\end{document}
